\section{原创创新点拆解}
\label{sec:innovation}

本研究从课程议题出发，经历了提出问题、实证否定与问题聚焦三个阶段。课上讲授的“海洋大数据的计算引擎”使作者联想到：航行器端侧的算力与能量有限，计算、测量与运动之间必须权衡。作者最初在东海赤潮巡测场景下研究机载计算与观测的时间分配，发现多算只会抬高规划器的自我估值（图\ref{fig:curse}），而时间分配本身的改进空间也很小：出发前模拟选出的分配策略平均遗憾为0.265，事后为每个任务挑选最优策略也只能降到0.243，远小于下水位置（约0.13）与信息新鲜度（约0.03）的影响。作者据此指出该方法增益小、新意不足，并把问题聚焦到选错一侧就要多花数小时的绕障选边：既然多算无法区分不同的真实海洋，能否直接检验是否存在与粗图相容、却使另一侧更快的合理反例？这一追问形成了本文的研究问题与表\ref{tab:innovation}所列的四个创新点。

\begin{table}[tb]
\centering
\caption{原创创新点、已有做法、实验证据与局限。}
\label{tab:innovation}
\footnotesize
\begin{tabular}{@{}p{0.13\linewidth}p{0.20\linewidth}p{0.22\linewidth}p{0.19\linewidth}p{0.14\linewidth}@{}}
\toprule
创新点 & 已有做法及其不足 & 本文做法 & 证据 & 局限 \\
\midrule
1. 区分数值歧义与环境歧义 & 加密网格或增加推演，把求解误差与信息缺失混为一谈。 & 以“能否找到合理反例”判定歧义的性质：数值歧义交给计算，环境歧义交给观测。 & 图\ref{fig:ambiguity}：加算64倍选错率仅从21.7\%降至19.7\%，换当天0.5°图降至10.4\%。 & 只在单个禁航区的左右选边上验证。 \\
\addlinespace
2. 用海洋大数据标定合理集合 & 凭经验设定不确定集，或借用历史相似日的差值。 & 用多年误差标定逐点误差、76千米相关长度的无散模态与覆盖90\%以上真实情况的半径。 & 图\ref{fig:data}、表\ref{tab:scope}：测试年份覆盖94\%；改用相似日差值后检验失效。 & 依赖历史统计，对台风等极端过程适应有限。 \\
\addlinespace
3. 最大后悔值的梯度搜索 & 随机抽样需上千个成员；改变选边的成员比例不计代价大小。 & 沿最快航线按包络定理求导，每次迭代2次求解，直接搜索最大后悔值。 & 图\ref{fig:risk}、\ref{fig:budget}：约71次求解达到2048次求解的效果；前20\%任务占53\%--71\%额外航时。 & 流场变化大时最快航线可能改变，一阶近似有误差。 \\
\addlinespace
4. 粗图不可信时先取新信息 & 不论风险一律观测或一律不观测。 & 最大后悔值超过阈值才付出观测代价；不能观测时改选最大后悔值较小的一侧。 & 图\ref{fig:main}、表\ref{tab:scope}：湾流与阿古拉斯回流区的额外航时降至1.31与2.20小时。 & 不观测时的保守改选只在黑潮主设置显著。 \\
\bottomrule
\end{tabular}
\end{table}

人机分工见附录B：研究问题与方向由作者提出并把关；Claude（Anthropic）完成文献检索、代码、实验与作图，并按作者要求重构撰写了本文终稿；此前各稿由Gemini（Google）根据Claude整理的事实清单撰写。
